% Generated draft: requires mathematical and source review.
\documentclass[11pt]{article}
\usepackage[margin=25mm]{geometry}
\usepackage{fontspec}
\setmainfont{Latin Modern Roman}
\usepackage{amsmath,amssymb,mathtools,unicode-math}
\setmathfont{Latin Modern Math}
\usepackage{enumitem,needspace,xurl,hyperref}
\setlength{\parindent}{0pt}
\setlength{\parskip}{5pt}
\setlength{\emergencystretch}{3em}
\newcommand{\N}{\mathbb{N}}
\newcommand{\Z}{\mathbb{Z}}
\newcommand{\Q}{\mathbb{Q}}
\newcommand{\R}{\mathbb{R}}
\newcommand{\C}{\mathbb{C}}
\newcommand{\F}{\mathbb{F}}
\newcommand{\A}{\mathbb{A}}
\newcommand{\PP}{\mathbb{P}}
\newcommand{\E}{\mathbb{E}}
\newcommand{\eps}{\varepsilon}
\newcommand{\dd}{\,\mathrm{d}}
\begin{document}
\section*{UnsolvedMath ID: 30004316}
\textbf{Source problem code:} catalogue entry
\subsection{Definitions and mathematical statement}
Let \(X\) be a set and let \(r:X\times X\to X\times X\) satisfy the braid form of the set-theoretic Yang--Baxter equation
\[
(r\times\mathrm{id}_X)(\mathrm{id}_X\times r)(r\times\mathrm{id}_X)
=
(\mathrm{id}_X\times r)(r\times\mathrm{id}_X)(\mathrm{id}_X\times r)
\quad\text{on }X^3.
\]
Here involutivity means \(r^2=\mathrm{id}_{X\times X}\), and idempotence means \(r^2=r\). Bijectivity of \(r\) is not imposed in this preliminary definition, since the supplied record explicitly discusses idempotent solutions; the precise intended class of solutions remains unspecified.

The supplied question is: \emph{Is the quantum symmetrizer \(\operatorname{QS}\) bijective for general set-theoretic Yang--Baxter solutions?} The record does not define this map, its domain, its codomain, or any coefficient or grading conventions. Consequently, it does not determine a complete mathematical statement of bijectivity.

To express only the logical form of the intended question, write
\[
\operatorname{QS}_{X,r}:D_{X,r}\longrightarrow C_{X,r},
\]
where \(D_{X,r}\) and \(C_{X,r}\) are placeholders for the unspecified source and target of the intended construction. Once that construction and the admissible class of solutions are fixed, the proposed assertion would be
\[
\text{for every admissible }(X,r),\qquad
\forall c\in C_{X,r}\;\exists!\,d\in D_{X,r}
\quad \operatorname{QS}_{X,r}(d)=c.
\]
This conditional formulation does not supply the missing construction or hypotheses. The supplied dated assessment reports that bijectivity is proved for involutive and idempotent solutions and remains open for general solutions. These are supplied literature claims, not independently verified mathematical conclusions.
\subsection{Short English statement}
The problem asks whether the quantum symmetrizer associated with every solution in the intended general class is bijective. A precise conjecture requires the definition of that map, its source and target, and the admissible hypotheses. The supplied record labels the general case open and reports positive results for involutive and idempotent solutions; these claims require independent review.
\subsection{Sources}
\begin{itemize}
\item Supplied local UnsolvedMath registry record, permanent ID \(30004316\), title \emph{Bijectivity of the Quantum Symmetrizer}, collection \emph{Oberwolfach Reports Open Problems}, exact source problem code \texttt{catalogue entry}. The record gives proposed year \(2019\) and supplies no individual proposer attribution.
\item Supplied report citation: \emph{Mini-Workshop: Algebraic Tools for Solving the Yang--Baxter Equation} (2020), Oberwolfach Reports, pp.\ 3207--3242, DOI \texttt{10.4171/owr/2019/51}. \url{https://doi.org/10.4171/owr/2019/51}. The supplied background describes the report as documenting workshop discussions and open problems.
\item Supplied literature source: \emph{Oberwolfach Reports, quantum symmetriser problem}. \url{https://ems.press/content/serial-article-files/46831}.
\item Supplied literature source: \emph{Cohomology of idempotent braidings}. \url{https://arxiv.org/abs/1607.08081}.
\item Supplied literature source: \emph{Plactic monoids: A braided approach}. \url{https://www.sciencedirect.com/science/article/am/pii/S002240492030418X}.
\item Literature assessment dated exactly \texttt{2026-08-22}: the supplied curation labels the problem open and reports bijectivity for involutive and idempotent solutions, with the general case remaining open. This is supplied dated curation derived from the OWR record, not an independently verified assessment; the background explicitly requests independent verification of time-sensitive status and cited claims.
\end{itemize}
\end{document}
